\section{Conclusion}\label{sec:conclusion}

We measured sixteen Philippine languages across three independent dimensions:
orthography, phonology, and geospatial features. Using trigram-based similarity
matrices and a purpose-built linguistic map. While the two quantitative
matrices agree on the ordering of language relationships, they diverge
significantly in scale, revealing two distinct historical layers. Orthographic
similarity ranges widely (0.3423 to 0.9391), registering colonial-era
standardization at full amplitude. It tracks recent, administratively tractable
histories of missionary, educational, and colonial contact. In contrast,
phonetic similarity operates within a substantially narrower band (0.7622 to
0.9927), preserving a compressed but legible record of shared Austronesian
descent that recent spelling reforms did not erase.

Geography accounts for these patterns where terrain restricted movement, but
fails where it did not; navigable water served to carry contact rather than
block it. Geographically, Romblomanon emerges as central under every measure,
while Northern Luzon forms a distinct zone bounded by the Cordillera and Sierra
Madre, neatly explaining the cohesion of the Bisayan cluster and Ilocano's
separation. Conversely, geography fails to explain Chavacano, which remains
universally peripheral due to creolization rather than physical isolation.
Similarly, geographic models poorly predict the Paranan-Pangasinan affinity,
which owes more to retained older features than to physical terrain.

Ultimately, similarity computed over trigrams cannot separate the retention of
conservative features from convergent innovation or undocumented contact.
Resolving these distinctions requires deeper historical evidence, including
cognate sets, reconstructed proto-forms, and migration histories. The languages
that most resist geographic explanation—foremost Yami and Paranan—are exactly
where this evidence would be most productively sought.
